\chapter{Goals}
\setcounter{section}{0}

\section{General Goal}

The general goal of this project is to develop a fully autonomous desktop automation system that enables users with limited mobility to use a Windows computer using only their voice, without requiring any manual intervention in the process.

This system is intended for users affected by conditions such as reduced motor impairment, due to their innability to operate traditional input devices, without learning a new interaction and without selling their data. Existing accessibility tools either demand step-by-step voice commands from the user, shifting the cognitive challenge back to them, or rely on cloud-based services that transmit sensitive data to external servers, where that data is used to train their own models, which may not be acceptable for users with privacy concerns.

This project addresses both limitations by delivering a system that accepts a single natural language voice prompt and executes the described task end-to-end on the local device, with no data leaving the machine at any point. The objective is not to build a faster or more convenient automation tool for technical users, but to ease the path of independent computer use for people who cannot rely on a keyboard and mouse, making any desktop genuinely accessible through natural speech.

\section{List of specific goals}

The following specific goals divide the general objective into concrete, measurable targets. Each goal addresses a distinct requirement that the system must satisfy to be considered a valid accessible automation solution.

\begin{enumerate}

    \item \textbf{Goal 1: Accept voice natural language prompts.} The user must be able to fully describe any task using only a spoken natural language prompt, with no keyboard, mouse, or touch input required from the user at any point. Additionally, the system must not impose a fixed command vocabulary or syntax, meaning it must be a free-form system where the same task can be expressible in different ways and still be correctly understood. This allows the user to interact with the system in a more flexible way, without needing to learn specific commands or phrasing, which is essential for users with cognitive or motor impairments.

    \item \textbf{Goal 2: Make the system fully autonomous.} Once the initial voice prompt has been received, the system must complete the task without any further human input. This autonomy extends beyond simple execution, the system must be capable of detecting when a step has not produced the expected result and act accordingly, without requiring the user to inspect the screen or provide feedback. A system that completes tasks but cannot tell whether it has succeeded is not a reliable accessibility tool.

    \item \textbf{Goal 3: Make the planning layer robust and reproducible.} The planning layer is perhaps the most important component, as it must generate correct, minimal action sequences for every task, with no mistaked, nor any redundancy. Therefore, the same or similar task descriptions must always produce the same correct plan, with no variation across runs. This is not just a quality requirement but a safety requirement, a user who cannot verify the screen or has severe issues to do so must be able to trust that the system always takes the shortest and most predictable path to the goal, without any unexpected detours or mistakes.

    \item \textbf{Goal 4: Achieve a task success rate of 90\%.} It is crucial that the system successfully completes every task it is designed to handle, with no partial executions or unresolved failures, making it precise and reliable.

    \item \textbf{Goal 5: Outperform other local and cloud-based competitors.} The system must not only achieve the highest task success rate among all benchmarked alternatives, both local and cloud-based, but it must do so without any trade-offs, making this system the clear choice on privacy and performance. It must use fewer steps, produce fewer errors, and expose no user data, all at the same time.

    \item \textbf{Goal 6: Complete every task within five minutes on consumer-grade hardware.} Every evaluation run must finish within 300 seconds from the moment the voice prompt is received, running on a standard consumer GPU without any hardware modifications or cloud offloading. A system that is theoretically accessible but takes an impractical amount of time to respond is not usable in practice. The five-minute limit is the threshold above which a user waiting for a result would reasonably abandon the interaction, try to intervene or look for an alternative solution.

    \item \textbf{Goal 7: Run the entire pipeline on local hardware.} The system must operate entirely on the user's own machine, with no external service, remote model, or cloud API involved at any point in the pipeline. All inference must run locally, making the system functional without an internet connection and independent of third-party availability, rate limits, or service outages. This ensures that the system remains usable in any environment, regardless of network access.

    \item \textbf{Goal 8: Guarantee full data privacy with zero data leaving the device.} No user data, including voice recordings, screenshots, task descriptions, file names, or intermediate outputs, are to be transmitted outside the device at any point during execution. The system may manage sensitive data and therefore, it must ensure that all information is stored and processed locally.

    \item \textbf{Goal 9: Interact with the desktop exclusively through standard graphical interface.} The system must operate the computer using only mouse clicks and keyboard on visible graphical elements, the same actions a human user would take. It must never open a terminal  or any other technical interface that would make a non-technical user uncomfortable. A system that solves tasks through the terminal may be functionally correct but minmises the trust of the user, who has no way to understand or verify what the hidden commands are doing.

\end{enumerate}

\section{Validation methods}

Each specific goal is validated through a concrete, measurable method applied during the experimental evaluation described in Chapter~\ref{chap:experiments}.

\begin{enumerate}

    \item \textbf{Goal 1: Accept voice natural language prompts.} All 21 evaluation runs are prompted through a free-form spoken natural language prompt, using Whisper. No keyboard or mouse interaction is performed by the evaluator during task execution.

    \item \textbf{Goal 2: Make the system fully autonomous.} The system organises the execution steps, thanks to Crew AI which helps orchestrate the entire workflow. In the system, the Manager Agent is responsible to orchestrate the entire process, placing each agent in charge of a specific atomic task, and depending on the output of each step to decide the next one. Finally, the Manager Agent, uses the QA Agent for self-verification, which is responsible to check if the task has been completed correctly, without any human feedback, additionally, all runs where evaluated by a human reviewer.
    
    \item \textbf{Goal 3: Make the planning layer robust and reproducible.} Each step across all 21 runs is inspected against the expected PDDL action sequence. Any wrong, missing, or redundant action is counted as a planning error. The validation criterion is \textit{M-04 Planning Error Rate}, which has to be 0. Reproducibility is validated by \textit{M-08 Standard Deviation of Steps}, which has to be 0, confirming that the same template always produces the same number of steps across all three runs.

    \item \textbf{Goal 4: Achieve a task success rate of 90\%.} The system is evaluated over 7 task templates with 3 runs each, for a total of 21 runs. A run is marked successful if the task is completed correctly as confirmed by a human review. The validation criterion is \textit{M-01 Task Success Rate}, which has to be 1 in at least 19 out of 21 runs.

    \item \textbf{Goal 5: Outperform other local and cloud-based competitors.} The system is benchmarked against five alternative systems in a common evaluation table covering all defined metrics. The validation criterion is that the proposed system achieves the best or joint-best value across M-01 TSR, M-04 Planning Error Rate, M-07 Total Error Rate, M-10 Data Exposure Level, and M-11 Total Steps simultaneously, with no metric where it is outperformed by any alternative.

    \item \textbf{Goal 6: Complete every task within five minutes on consumer-grade hardware.} End-to-end latency is measured for each of the 21 runs from the moment the prompt is received to task completion. The validation criterion is \textit{M-03 End-to-End Latency}, which has to be below 300 seconds for every individual run, with no exceptions.

    \item \textbf{Goal 7: Run the entire pipeline on local hardware.} The system is validated by confirming with each run that WiFi and Ethernet connections are disabled, ensuring that no network access is available. Additionally, the system uses all free, local and open-source tools and models.

    \item \textbf{Goal 8: Guarantee full data privacy with zero data leaving the device.} Privacy is validated with each run by ensuring that WiFi and Ethernet connections are disabled, confirming that no data leaves the device, not even after the execution is done. This is reported as \textit{M-10 Data Exposure Level}, where:
        \begin{itemize}
            \item 0: No data leaves the device.
            \item 1: Metadata or task text sent externally.
            \item 2: Sensitive data sent externally.
            \item 3: Sensitive data used to train models.
        \end{itemize}
        The validation criterion is a score of 0.

    \item \textbf{Goal 9: Interact with the desktop exclusively through standard graphical interface.} There is evidence for each run that the system only uses mouse clicks and keyboard inputs on visible graphical elements, with no terminal commands or other technical interfaces involved. This is confirmed by the logs automatically generated, alongside the screenshots taken throughout the process and a final human reviewer.

\end{enumerate}
